\subsubsection{SEM Usage}\label{SEC:SEM:usage}

Having described some of the basic theory and nomenclature of the
synthetic eddy method, we will now discuss how it works within the
GFR code.

%###############################################################################

\subsubsubsection{Splitting off Subsonic BC for Boundary Faces
Not in SEM Profile}[SEC:SEM:usage:subsonic]

In general, an SEM boundary condition is normally specified for the entire
inflow boundary, but the height of the boundary layer defined in the SEM profile
is usually only a small fraction of the overall inflow.
%
Since the computational cost of the SEM scales with the number of boundary
points involved in the SEM, any boundary points that are outside the height of
the SEM profile, \hSEM, are split off into a new and completely separate
subsonic inflow boundary condition that uses the standard method for a subsonic
inflow.
%
In order to accomplish this, during the pre-processor, the wall distances for
all boundary points on the SEM boundary are computed and then sorted to separate
those boundary points with wall distance exceeding $\hSEM$ into the new subsonic
inflow boundary condition.
%
The flow conditions for this new subsonic inflow boundary are a direct copy of
the flow conditions specified in
\BCinpBC[nmlBC::BCInput][i]{}
for the respective SEM boundary condition.
%
This ensures that 1) the minimum number of boundary points are contained within
the SEM boundary condition, and 2) the subsonic inflow outside the SEM is not
overprescribed allowing for characteristic information to both enter and leave
the computational domain through the boundary.

%%%%%%%%%%%%%%%%%%%%%%%%%%%%%%%%%%%%%%%%%%%%%%%%%%%%%%%%%%%%%%%%%%%%%%%%%%%%%%%%

There are two details to keep in mind regarding this behavior of splitting off
boundary points not within the SEM profile height into a new separate boundary
condition.
%
\begin{itemize}
  %
  \item \textit{Consistency between SEM and subsonic inflow boundary
    conditions.}

    It is best to make sure that the flow conditions specified in
    \BCinpBC[nmlBC::BCInput][i]{}
    for an SEM boundary condition are consistent with the flow
    conditions at the boundary layer edge of the SEM profile file.
    %
    Stability issues can occur if there is a mismatch between the SEM profile
    and the corresponding subsonic inflow boundary condition.
    %
    Since the subsonic inflow boundary conditions generally prescribe a fixed
    total temperature or total pressure, it is recommended to check that these
    two quantities are consistent between the specified boundary condition
    inputs and their respective values at the outer edge of the SEM profile.
    %
    Other possible quantities to check are Mach number and velocity magnitude.
  %
  \item \textit{SEM profile height and efficiency.}

    The profile data in the SEM profile file is not required to only contain the
    boundary layer part of the flow.
    %
    The profile is allowed to contain as much of the main freestream flow as
    desired.
    %
    There can be some inefficiencies from including part of the freestream flow
    in the SEM profile.
    %
    The height of the SEM box is based on the height of the SEM profile, \hSEM;
    and the number of synthetic eddies is based on the volume of the SEM box,
    \key[key:SEMVolume]{\vSEM}[Synthetic Eddy Method (SEM)!volume].
   %\vSEM\key[key:SEMVolume]{\vSEM}[Synthetic Eddy Method (SEM)!volume].
    %
    Therefore, including parts of the freestream flow will increase the
    number of synthetic eddies in the SEM box.
    %
    In the freestream, the turbulent stresses are generally near-zero.
    %
    When the SEM algorithm computes a random fluctuation for a boundary point in
    the freestream section, the correlation converting this to a physical
    fluctuation will result in a non-fluctuating velocity profile since the
    turbulence is zero.
    %
    All of this is to say that the freestream parts of the SEM profile will
    waste computational resources trying to compute velocity fluctuations where
    there are no fluctuations.

    It is generally best to only specify the boundary layer profile up to, but
    not beyond, the boundary layer edge, \thatis $\hSEM = \lBC$.
    %
    By specifying the minimum profile height, more boundary points will get
    separated to the new subsonic inflow boundary condition.
    %
    The traditional subsonic inflow BC is more efficient than the SEM BC.
    %
    Even with the highly efficient manner in which the SEM BC has been
    implemented, the traditional subsonic inflow BC is still more efficient; the
    difference between a full subsonic inflow BC and a full SEM inflow BC is
    approximately around 10-15\% additional computational cost.
  %
\end{itemize}

%###############################################################################

\subsubsubsection{Restarting Simulations with SEM}[SEC:SEM:usage:restart]

When using one or more SEM boundary conditions, the state of each SEM is
written to the restart file, where it is appended to the end after all the
standard restart information is written.
%
The SEM information written to the restart file includes:
%
\begin{itemize}
  %
  \item The random seed information to allow continuing the random number
    generator sequence from exactly where it left off the previous time step.
  %
  \item The number of SEM boxes and the number of eddies in each box.
  %
  \item For each individual eddy: its location, its random intensities in each
    direction, and its length scale in each direction.
  %
\end{itemize}

%%%%%%%%%%%%%%%%%%%%%%%%%%%%%%%%%%%%%%%%%%%%%%%%%%%%%%%%%%%%%%%%%%%%%%%%%%%%%%%%

\noindent
There are three situations that can occur involving simulation restarts and SEM
boundary conditions:
%
\begin{itemize}
  %
  \item \textit{Restarting a simulation that previously used an SEM boundary
    condition and continuing use of the SEM.}
     
    This is probably the most common of the three situations.
    %
    The SEM information will be read from the restart file and a few sanity
    checks will be made to make sure the restart SEM is consistent with what is
    given in the input file.
    %
    For example, the location of each eddy will be checked against the SEM boxes
    created from the information in the input file, and execution will be
    stopped with an error message if any eddy is located outside the SEM boxes.

    This is also the only scenario in which the input variable
    \inp[nmlSEM::UseRandomSeedFromRestart]{Use\_Random\_Seed\_From\_Restart}
    has any effect on a simulation.
    %
    This variable controls whether or not the random seed information in the
    restart file is used when initializing the SEM algorithm at the start of
    the new\slash restarted simulation.
    %
    If set to \tru, the random seed from the restart will be used to reset the
    random number generator to exactly where it left off when the restart file
    was written, and the SEM will behave as if the simulation never stopped.
    %
    The default value of \fls{} will ignore the random seed in the restart file,
    and it will instead use the current date and time to create a new
    \quoterm{random} seed for the random number generator.
    %
    This might have a very small effect on the \quoterm{randomness} of the
    locations and intensities generated for new eddies needed to replace eddies
    that have propagating completely through their SEM box.
    %
    However, the total number of time steps needed for a simulation will almost
    certainly prevent this from having any kind of measurable effect on the
    simulation results.

    NOTE: Very rarely (but it does happen), the SEM can randomly cause a
    simulation to abort execution due to an NaN or other floating point
    exception.
    %
    Usually, the only way to get past one of these occasions without modifying
    the solution through limiting \andor filtering is to reset the random
    number generator down a different path by setting
    \\ % Prevent next word from entering right margin.
    \inp[nmlSEM::UseRandomSeedFromRestart]{Use\_Random\_Seed\_From\_Restart}%
    \texttt{=}\fls{}.
  %
  \item \textit{Restarting a simulation that previously used an SEM boundary
    condition and replacing it with a non-SEM boundary condition.}
     
    This scenario is particularly useful when replacing the SEM with a periodic
    boundary condition.
    %
    An example of this would be something like a turbulent channel flow case.
    %
    Starting a streamwise periodic turbulent channel flow from a uniform
    solution would take an enormous amount of simulation (and computational)
    time to develop into a fully turbulent flow without any sort of tripping
    mechanism (which you would not want for a streamwise periodic flow).
    %
    One way to speed this up is to use an SEM inflow with a subsonic outflow
    boundary condition and run the simulation for just over one
    flow-through-time to populate the channel with turbulent structures.
    %
    Then the simulation could be restarted replacing the SEM inflow and
    subsonic-outflow boundary conditions with a streamwise periodic
    boundary condition.
    %
    In this scenario, the SEM information is read from the restart file but it
    will be completely ignored since there are no SEM boundary conditions
    specified in the input file.\\
  %
  \item \textit{Adding an SEM boundary condition when restarting a simulation
    that did not use one.}
     
    This scenario might occur when running a simulation without an SEM to get a
    stable flowfield as a pseudo-initialization and then add an SEM afterwards,
    or perhaps to add an SEM to the inflow to see how that might
    affect the overall solution.
    %
    In this scenario, the SEM is basically set up in the exact same way as if
    you were starting a new simulation, with the difference being you maintain
    whatever time history information exists in the restart file.
  %
\end{itemize}

%###############################################################################

\subsubsubsection{SEM Profile File}[SEC:SEM:usage:profile]

The following three sections describe: the required format of the SEM profile
file (\ref{SEC:SEM:usage:profile:format}); the valid variables, and their
respective identifier strings, that can be included in the SEM profile file
(\ref{SEC:SEM:usage:profile:variables}); and how the SEM profile data is
mapped onto the boundary face points using interpolation and a rotation
transformation (\ref{SEC:SEM:usage:profile:mapping}).

%###############################################################################

\FourSubsection{SEM Profile File Format}[SEC:SEM:usage:profile:format]

% Command to format data values (DVs) of SEM profile file.
%   \DV{i}{j} with i=1:n, j=1:Np. i=variable number, j=data point.
\newcommand{\DV}[2]{V$_{#1}(#2)$}

\begin{table}[h!]
  \texttt{
    \begin{tabular}{lllll}
      \toprule
      "V$_1$"     & "V$_2$"     & "V$_3$"     & \ldots{} & "V$_n$"     \\
      Np          &             &             &          &             \\
      \DV{1}{1}   & \DV{2}{1}   & \DV{3}{1}   & \ldots{} & \DV{n}{1}   \\
      \DV{1}{2}   & \DV{2}{2}   & \DV{3}{2}   & \ldots{} & \DV{n}{2}   \\
      \ldots{}    &             &             &          &             \\
      \DV{1}{\texttt{Np}} &
      \DV{2}{\texttt{Np}} &
      \DV{3}{\texttt{Np}} &
      \ldots{}            &
      \DV{n}{\texttt{Np}} \\
      \bottomrule
    \end{tabular}
  }
\end{table}

%%%%%%%%%%%%%%%%%%%%%%%%%%%%%%%%%%%%%%%%%%%%%%%%%%%%%%%%%%%%%%%%%%%%%%%%%%%%%%%

In the above file format:% \\
%
\begin{itemize}
  \item
  Columns are separated by space(s).
  Multiple spaces are treated as one.
  Leading and trailing spaces on each line are ignored.
%
  \item
  $n$ is the number of variables present in the file.
  It is also the number of columns in the file.
%
  \item
  \texttt{Np} is an integer giving the number of profile data locations,
  \thatis the number of rows of data.
%
  \item
  \texttt{"V$_i$"} is a character string specifying the variable
  name for the $i$\supth{} column of data.
  Quotes are suggested but not required.
  Valid names are described in \sref{SEC:SEM:usage:profile:variables}.
%
  \item
  \texttt{\DV{i}{j}} is a floating point value for the $j$\supth{}
  profile point of the $i$\supth{} variable, where
  $j=1,\ldots,\mathtt{Np}$ and $i=1,\ldots,n$.
  Data values are specified in SI units.
  %
  The wall distance variable, $d$, should be in
  \key[key:GridScaledUnits]{grid scaled units}[grid scaled units]
  (\thatis meters).
  The profile will be imposed on the mesh points
  (after the \inp[nmlGrid::GridScalingFactor]{Grid\_Scaling\_Factor}
  has been applied to the mesh).
%
% \item
% \texttt{\DV{i}{j}} data values are specified in SI units.
%
  \item
  The data rows should be sorted in the order of increasing
  distance from the wall.
%
\end{itemize}

%###############################################################################

\FourSubsection{SEM Profile Variable Names}[SEC:SEM:usage:profile:variables]
 
\vspace{0.5em}
%\rule{0.75in}{0.25em} tabcolsep=\the\tabcolsep\\
%
\begingroup % Contain the scope of \growsep.
\newlength{\growsep} % extra space for group row separation
\setlength{\growsep}{0.50em}
% Default \tabcolsep = 6pt, so 1 column gap = 12pt, and 2=24pt=0.33in.
% Manually adjust col3 width to match \rule{\textwidth}{0.1in}.
%\begin{tabular}{p{0.75in} p{2.0in} p{3.25in}}% 6.5-0.75-2.0-0.33=3.42
%\begin{tabular}{c p{2.0in} p{3.25in}}% 6.5-0.75-2.0-0.33=3.42
\begin{tabular}{c @{\qquad} l @{\qquad} p{3.25in}}% 6.5-0.75-2.0-0.33=3.42
  Variable    & Input Name & Description \\
  \hline
  $d$         & \texttt{Y R}
              & wall distance \\
  $\rho$      & \texttt{RHO DENSITY}
              & density \\
  $u$         & \texttt{U VX}
              & axial velocity, aligned with the mean flow of the profile \\
  $v$         & \texttt{V VY}
              & vertical velocity, perpendicular to the wall \\
  $w$         & \texttt{W VZ}
              & span velocity, parallel to the wall and perpendicular to the
                mean flow \\
  $T$         & \texttt{T TEMPERATURE}
              & temperature \\[\growsep]
  $\tau_{xx}$ & \texttt{TAUXX TAU11 RXX R11 UU}
              & $\ReyAve{\ReyFlx{u}\ReyFlx{u}}$ normal stress
                \\[\growsep]
  $\tau_{yy}$ & \texttt{TAUYY TAU22 RYY R22 VV}
              & $\ReyAve{\ReyFlx{v}\ReyFlx{v}}$ normal stress
                \\[\growsep]
  $\tau_{zz}$ & \texttt{TAUZZ TAU33 RZZ R33 WW}
              & $\ReyAve{\ReyFlx{w}\ReyFlx{w}}$ normal stress
                \\[\growsep]
  $\tau_{xy}$ & \texttt{TAUXY TAU12 RXY R12 UV}
              & $\ReyAve{\ReyFlx{u}\ReyFlx{v}}$ shear stress ($=\tau_{yx}$) \\
              & \texttt{TAUYX TAU21 RYX R21 VU}
              & \\[\growsep]
  $\tau_{xz}$ & \texttt{TAUXZ TAU13 RXZ R13 UW}
              & $\ReyAve{\ReyFlx{u}\ReyFlx{w}}$ shear stress ($=\tau_{zx}$) \\
              & \texttt{TAUZX TAU31 RZX R21 WU}
              & \\[\growsep]
  $\tau_{yz}$ & \texttt{TAUYZ TAU23 RYZ R23 VW}
              & $\ReyAve{\ReyFlx{v}\ReyFlx{w}}$ shear stress ($=\tau_{zy}$) \\
              & \texttt{TAUZY TAU32 RZY R32 WV}
              & \\
  \hline
\end{tabular}
\endgroup
%\\ \rule{\textwidth}{0.10em}%  For visually sizing the table width.
%\\[1.0em]
%\FIXME{
%\begin{itemize}
%\item
%Is $u$ actually in the direction of the (inward-pointing) inflow plane
%surface normal?
%\item
%For a backwards channel flow moving right-to-left, which direction does
%$u$ point?
%Is $u$ positive?
%Is \BCinptext[i]{SEM\%Convective\_Velocity} positive?
%\item
%Is $v$ in the wall-normal direction or in the direction of the
%wall-normal projected onto the inflow plane?
%(eg, if the inflow plane is in the convergent section of a nozzle,
%is $v$ radial? Otherwise, how is it normal to $u$?) 
%\end{itemize}
%}

\vspace{1.0em}%
\noindent
A few things to note about these variables:
\begin{itemize}
  \item These strings are case insensitive when read in by the SEM code.
  \item The wall distance variable, $d$, is the only required quantity in the
    SEM profile file. All other variables are optional.
    \begin{itemize}
      \item If any of the mean flow variables ($\rho, u, v, w, T$) are not
        provided, the value of that variable in the profile will be set to a
        constant based on the conditions specified for the respective BC in
        \BCinpBC[nmlBC::BCInput][i]{}.
        It should be noted that it is probably best that the
        axial velocity, $u$, also be specified since the intended use is for a
        boundary layer.
      \item If any of the six stress terms are not provided, the value of that
        variable will be set to zero for the entire profile.
    \end{itemize}
  \item The wall distance variable, $d$, should be in
    \key[key:GridScaledUnits]{grid scaled units}[grid scaled units]
    (\thatis meters).
    The profile will be imposed on the mesh points 
    (after the \inp[nmlGrid::GridScalingFactor]{Grid\_Scaling\_Factor}
    has been applied to the mesh).
  \item The stress variables in the profile file are treated as
    $\tau_{ij} = \ReyAve{\ReyFlx{u}_{i}\ReyFlx{u}_{j}}$
    in units of \unit{\m\squared\per\s\squared},
    without the density or minus sign often used in the definition of
    turbulent stress.
    %
    These would more correctly be referred to as single-point velocity
    correlation terms.
  \item Using the string \quoterm{\texttt{R}} for wall distance tells the code
    to use a cylindrical box instead of a rectangular box. This has the same
    effect as adding
    \BCinpBC[nmlBC::BCInput/SEM/IsCylindrical][i]{SEM\%Is\_Cylindrical}%
    \texttt{=}\tru{}
    to the input file, where $i$ is the index for the SEM BC of interest.
    This solely affects how the eddies are randomly generated to achieve a
    truly random distribution in the radial, $r$, direction.
    %
    For a cylindrical SEM box, randomly generating
    $r\ (r_{\textrm{min}}\le r \le r_{\textrm{max}})$ and
    $\theta\ (\theta_{\textrm{min}}\le \theta \le \theta_{\textrm{max}})$
    coordinate values to place a new eddy will result in
    more eddies near the symmetry axis and
    less eddies further away.
    %
    In order to generate a truly uniform random distribution,
    we need to use the equation
      $r=\sqrt{\mathtt{Random}()\left[ r^2_{\textrm{max}}
                                     - r^2_{\textrm{min}}
                                \right] + r^2_{\textrm{min}}}$
    where $\mathtt{Random}()$ is a function that generates a random number
    on the unit interval.
    %
    This applies an area weighting to the radial distribution and
    results in a uniform distribution of eddies in the volume.

\end{itemize}

%###############################################################################

\FourSubsection{Mapping the SEM Profile onto Boundary Points}%
[SEC:SEM:usage:profile:mapping]

From the description of the SEM algorithm in sections
\ref{SEC:SEM:theory:DefineProfile} and \ref{SEC:SEM:theory:ComputeBC}, each
boundary face point included in the SEM inflow must have a mean flow state
\(\overline{\cvVec}\) and Reynolds stress tensor $R_{ij}$ prescribed from the
SEM profile data.
%
From \sref{SEC:SEM:usage:profile:variables}, the profile data is defined as a 1D
function of wall distance within a `wall-aligned coordinate system'.
%
In order to get the prescribed conditions at a boundary point from this 1D data,
two operations need to be performed on the data: interpolation and rotation.

%%%%%%%%%%%%%%%%%%%%%%%%%%%%%%%%%%%%%%%%%%%%%%%%%%%%%%%%%%%%%%%%%%%%%%%%%%%%%%%%

Interpolation is needed because the boundary points are almost certainly not
located at the exact wall distances given in the SEM profile, so the first step
of interpolation is to compute the wall distance (\thatis distance from a given
point to the boundary point closest to it with a wall boundary condition) for
all SEM boundary points.
%
Once the wall distance for each SEM boundary point is known,
then the boundary layer properties can be obtained using a simple linear
interpolation of the provided SEM profile file.

%%%%%%%%%%%%%%%%%%%%%%%%%%%%%%%%%%%%%%%%%%%%%%%%%%%%%%%%%%%%%%%%%%%%%%%%%%%%%%%%

Because the profile data is given in a `wall-aligned coordinate system', a
rotation also needs to be applied to the velocity vector and stress tensor
parts of the
interpolated profile data such that it is aligned relative to how the wall
normal of the nearest wall point is orientated within the Cartesian coordinate
system.
%
The remainder of this section consists of a list of three bullet points
providing details on how the wall normal vector is used to create the
transformation matrix needed for rotating the SEM profile data defined in a
`wall-aligned coordinate system' to the standard Cartesian coordinate system of
the computational domain.
%
This information is included here mostly for reference \andor code development
purposes.
%
Therefore, the general user only interested in how to specify and use an
SEM boundary condition in a simulation can skip the remainder of this section.

%%%%%%%%%%%%%%%%%%%%%%%%%%%%%%%%%%%%%%%%%%%%%%%%%%%%%%%%%%%%%%%%%%%%%%%%%%%%%%%%

\begin{itemize}[parsep=2.5\parsep]

%-------------------------------------------------------------------------------

  \item \textit{Coordinate Transforms Using a Wall-Aligned Normal Vector}

    In the description that follows, we are interested in mapping the SEM
    profile data onto the boundary points on the SEM inflow plane.
    %
    The SEM data is provided as a 1D profile relative to wall distance and
    extending in the direction of the wall-normal vector.
    %
    Therefore, our goal is to create a transformation matrix, \rmat, that
    transforms the SEM vector\slash tensor data from the `wall-aligned
    coordinate system' to the standard Cartesian coordinate system of the grid.

%%%%%%%%%%%%%%%%%%%%%%%%%%%%%%%%%%%%%%%%%%%%%%%%%%%%%%%%%%%%%%%%%%%%%%%%%%%%%%%%

    Let \((\xsb,\ysb,\zsb)\) be the standard unit basis vectors for the
    Cartesian coordinate system in \(\mathbb{R}^3\), and let
    \((\xwb,\ywb,\zwb)\) be the unit basis vectors for the wall-aligned
    coordinate system also in \(\mathbb{R}^3\).
    %
    With \(\wac{v}=v_{\wac{x}}\xwb+v_{\wac{y}}\ywb+v_{\wac{z}}\zwb\) being an
    arbitrary vector defined by the wall-aligned basis vectors, our goal is to
    find \rmat{} such that \[\vec{v}=\rmat\cdot\wac{v}\] where
    \(\vec{v}=v_x\xsb+v_y\ysb+v_z\zsb\) is the same vector but defined by the
    standard Cartesian basis vectors.\\
    %
    NOTE:
    Our definition of \rmat{} may be the inverse of what is found in other
    sources, \thatis they might define it such that
    $\wac{v}=\bm{\breve{\rvar}}\cdot\vec{v}$.
    %
    Since one of the properties of transformation matrices is
    $\rmat^{\mathsf{T}}=\rmat^{-1}$,
    this just means our \rmat{} is equivalent to their
    $\bm{\breve{\rvar}}^{\mathsf{T}}$.
    
%%%%%%%%%%%%%%%%%%%%%%%%%%%%%%%%%%%%%%%%%%%%%%%%%%%%%%%%%%%%%%%%%%%%%%%%%%%%%%%%

    Before we continue to define \rmat, we first need to define two surface
    normals for the relevant grid boundaries.
    %
    Let \[\normal=\nx\xsb+\ny\ysb+\nz\zsb\] be the unit vector that is normal to
    the wall surface and points out of the wall and into the computational
    domain.
    %
    Similarly, let \[\innorm=\hat{s}_x\xsb+\hat{s}_y\ysb+\hat{s}_z\zsb\] be a
    unit vector that represents some sort of \quoterm{collective unit-normal}
    for the SEM inflow plane that points into the computational domain.
    %
    A unit vector based on the direction of the inflow velocity vector could
    also be used here.
    
%%%%%%%%%%%%%%%%%%%%%%%%%%%%%%%%%%%%%%%%%%%%%%%%%%%%%%%%%%%%%%%%%%%%%%%%%%%%%%%%

    The general method for formulating a transformation matrix is
    %
    \begin{equation*}
      \rmat = 
      \begin{bmatrix}
        \cos(\xwb,\xsb) & \cos(\ywb,\xsb) & \cos(\zwb,\xsb)\\
        \cos(\xwb,\ysb) & \cos(\ywb,\ysb) & \cos(\zwb,\ysb)\\
        \cos(\xwb,\zsb) & \cos(\ywb,\zsb) & \cos(\zwb,\zsb)
      \end{bmatrix}
    \end{equation*}
    %
    where for any two vectors \(\vec{u}\) and \(\vec{v}\)
    %
    \begin{equation*}
      \cos\left(\vec{u},\vec{v}\right)=
          \frac{\vec{u}\cdot\vec{v}}{\|\vec{u}\| \|\vec{v}\|}
    \end{equation*}
    %
    One way to think of \rmat{} is that each column is a basis vector of the
    wall-aligned coordinate system expressed as a unit vector in the standard
    Cartesian coordinate system.
    %
    We can then express this slightly differently as
    %
    \begin{equation*}
      \rmat = 
      \begin{bmatrix}
        a_x & b_x & c_x \\
        a_y & b_y & c_y \\
        a_z & b_z & c_z
      \end{bmatrix}
    \end{equation*}
    %
    where we define the following three unit vectors as
    %
    \begin{alignat*}{5}
      \hat{a} & = \xwb & = a_x\xsb & + a_y\ysb & + a_z\zsb \\
      \hat{b} & = \ywb & = b_x\xsb & + b_y\ysb & + b_z\zsb \\
      \hat{c} & = \zwb & = c_x\xsb & + c_y\ysb & + c_z\zsb
    \end{alignat*}

%%%%%%%%%%%%%%%%%%%%%%%%%%%%%%%%%%%%%%%%%%%%%%%%%%%%%%%%%%%%%%%%%%%%%%%%%%%%%%%%

    We can now start filling this transformation matrix.
    %
    The wall-aligned coordinate system should be just that, so we define the
    \ywb{} basis vector to be aligned with the wall normal which is simply
    \(\hat{b}=\ywb=\normal\) or \((b_x,b_y,b_z)=(\nx,\ny,\nz)\).
    %
    We know \(\hat{c}\) is perpendicular to \(\hat{b}\), so we just need to take
    the cross-product of some arbitrary vector with \(\hat{b}\) to satisfy this
    requirement.
    %
    Using \innorm{} as this arbitrary vector, we can define \(\hat{c}\) as
    %
    \begin{equation*}
      \begin{split}
        \hat{c} & = \frac{\innorm\times\hat{b}}{\|\innorm\times\hat{b}\|} \\
                & = \frac{\innorm\times\normal}{\|\innorm\times\normal\|}
      \end{split}
    \end{equation*}
    %
    Finally, knowing that the vector \(\hat{a}\) must be perpendicular to both
    \(\hat{b}\) and \(\hat{c}\), we can define \(\hat{a}\) as
    %
    \begin{equation*}
      \begin{split}
        \hat{a} & = \frac{\hat{b}\times\hat{c}}{\|\hat{b}\times\hat{c}\|} \\
                & = \frac{\normal\times(\innorm\times\normal)}%
                         {\|\normal\times(\innorm\times\normal)\|}
      \end{split}
    \end{equation*}
    %
    This can also be written in a slightly different way using the vector triple
    product identity
    %
    \begin{equation*}
      \hat{a} = \frac{\innorm - \left(\normal\cdot\innorm\right)\normal}
                     {\|\innorm - \left(\normal\cdot\innorm\right)\normal\|}
    \end{equation*}
    %
    Now that the three vectors \([\hat{a},\hat{b},\hat{c}]\) are defined, we can
    directly compute the matrix components of \rmat, which are now entirely in
    terms of the two known vectors \normal{} and \innorm.

%%%%%%%%%%%%%%%%%%%%%%%%%%%%%%%%%%%%%%%%%%%%%%%%%%%%%%%%%%%%%%%%%%%%%%%%%%%%%%%%

    As an example, we will assume the inflow plane is constant \(x\) resulting
    in \(\innorm=(1,0,0)\).
    %
    After some simplifying, the resulting transformation matrix is
    %
    \begin{equation*}
      \rmat =
      \begin{bmatrix}
        c                 & \nx & 0 \\[1.0ex]
        \frac{-\nx\ny}{c} & \ny & \frac{-\nz}{c} \\[1.0ex]
        \frac{-\nx\nz}{c} & \nz & \frac{ \ny}{c}
      \end{bmatrix}
    \end{equation*}
    %
    where \(c = \sqrt{ \ny^2 + \nz^2 }\).
    
%-------------------------------------------------------------------------------

\item \textit{Coordinate Transformation of a Vector}

    As stated previously, our goal is to use this transformation matrix to map
    the velocity vector and Reynolds stress tensor in the SEM profile file onto
    the SEM inflow boundary points.
    %
    For example, say we have a velocity vector \((\wac{u},\wac{v},\wac{w})\)
    from the SEM profile defined in the wall-aligned coordinate system.
    %
    In the case of a constant $x$ inflow plane,
    transforming this vector to the standard Cartesian coordinate system is
    simply
    %
    \begin{equation*}
      \begin{bmatrix}
        u \\ v \\ w
      \end{bmatrix} =
      \begin{bmatrix}
        c                 & \nx & 0 \\[1.0ex]
        \frac{-\nx\ny}{c} & \ny & \frac{-\nz}{c} \\[1.0ex]
        \frac{-\nx\nz}{c} & \nz & \frac{ \ny}{c}
      \end{bmatrix}
      \begin{bmatrix}
        \wac{u} \\ \wac{v} \\ \wac{w}
      \end{bmatrix}
    \end{equation*}
    %
    Expanding this out gives
    %
    \begin{align*}
      u & =        \wac{u}        c +\wac{v}\nx \\[1.0ex]
      v & = \frac{-\wac{u}\nx\ny}{c}+\wac{v}\ny-\frac{\wac{w}\nz}{c} \\[1.0ex]
      w & = \frac{-\wac{u}\nx\nz}{c}+\wac{v}\nz+\frac{\wac{w}\ny}{c}
    \end{align*}

%-------------------------------------------------------------------------------

\item \textit{Coordinate Transformation of a Tensor}

    Transforming a tensor in wall-aligned coordinates, \(\wac{\bm{\sigma}}\), to
    reference coordinates, \(\bm{\sigma}\), is accomplished through
    %
    \begin{equation*}
      \bm{\sigma} = \rmat \cdot \wac{\bm{\sigma}} \cdot \rmat^{\mathsf{T}}
    \end{equation*}
    %
    If the tensor \(\wac{\bm{\sigma}}\) is symmetric, as in the case of the
    Reynolds stress tensor, we are only concerned about the terms
    \([\sigma_{xx}, \sigma_{yy}, \sigma_{zz}, \sigma_{xy}, \sigma_{xz},
    \sigma_{yz}]\).
    %
    For a constant $x$ inflow plane,
    expanding out the above equation and solving for these terms results in a
    set of six linear equations that can be rewritten in matrix form as
    %
    \begin{equation*}
      %%%%%%%%%%%%%%%%%%%%%%%%%%%%%%%%%%%%%%%%%%%%%%%%%%%%%%%%%%%%%%%%%%%%%%%%%%
      \begin{bmatrix}
        \sigma_{xx}\\ \sigma_{yy}\\ \sigma_{zz}\\
        \sigma_{xy}\\ \sigma_{xz}\\ \sigma_{yz}
      \end{bmatrix} =
      %%%%%%%%%%%%%%%%%%%%%%%%%%%%%%%%%%%%%%%%%%%%%%%%%%%%%%%%%%%%%%%%%%%%%%%%%%
      \begin{bmatrix}
        %
        % First row for \sigma_{xx} --------------------------------------------
        c^2 & \nx^2 & 0 & 2c\nx & 0 & 0 \\[1.0ex]
        % Second row for \sigma_{yy} -------------------------------------------
        \frac{\nx^2\ny^2}{c^2} &
        \ny^2 &
        \frac{\nz^2}{c^2} &
        \frac{-2\nx\ny^2}{c} &
        \frac{2\nx\ny\nz}{c^2} &
        \frac{-2\ny\nz}{c} \\[1.0ex]
        % Third row for \sigma_{zz} --------------------------------------------
        \frac{\nx^2\nz^2}{c^2} &
        \nz^2 &
        \frac{\ny^2}{c^2} &
        \frac{-2\nx\nz^2}{c} &
        \frac{-2\nx\ny\nz}{c^2} &
        \frac{2\ny\nz}{c} \\[1.0ex]
        % Fourth row for \sigma_{xy} -------------------------------------------
        -\nx\ny &
        \nx\ny &
        0 &
        c\ny-\frac{\nx^2\ny}{c} &
        -\nz &
        \frac{-\nx\nz}{c} \\[1.0ex]
        % Fifth row for \sigma_{xz} --------------------------------------------
        -\nx\nz &
        \nx\nz &
        0 &
        c\nz-\frac{\nx^2\nz}{c} &
        \ny &
        \frac{\nx\ny}{c} \\[1.0ex]
        % Sixth row for \sigma_{yz} --------------------------------------------
        \frac{\nx^2\ny\nz}{c^2} &
        \ny\nz &
        \frac{-\ny\nz}{c^2} &
        \frac{-2\nx\ny\nz}{c} &
        \frac{\nx\left(\nz^2-\ny^2\right)} {c^2} &
        \frac{\ny^2-\nz^2}{c}
        %
      \end{bmatrix}
      %%%%%%%%%%%%%%%%%%%%%%%%%%%%%%%%%%%%%%%%%%%%%%%%%%%%%%%%%%%%%%%%%%%%%%%%%%
      \begin{bmatrix}
        \wac{\sigma}_{xx}\\ \wac{\sigma}_{yy}\\ \wac{\sigma}_{zz}\\
        \wac{\sigma}_{xy}\\ \wac{\sigma}_{xz}\\ \wac{\sigma}_{yz}
      \end{bmatrix}
      %%%%%%%%%%%%%%%%%%%%%%%%%%%%%%%%%%%%%%%%%%%%%%%%%%%%%%%%%%%%%%%%%%%%%%%%%%
    \end{equation*}

%-------------------------------------------------------------------------------

\end{itemize}

%###############################################################################

\subsubsubsection{Inputs Controlling Global SEM Behavior}%
[SEC:SEM:usage:otherinputs]

\begin{nmldescitem}
&nml_sem
  Use_Profile_For_Initialization = F
  Use_Random_Seed_From_Restart   = F
  Use_Spherical_Eddy_Volume      = F
  Include_Thermo_Fluctuations    = F
  Box_Num_Eddies_Multiplier      = 1.0
/
\end{nmldescitem}

Inputs that are universally applied to all SEM boundary conditions are
specified in the \inp[nmlSEM]{nml\_sem} input namelist block.
Each one controls a different part of the SEM algorithm.
%
\begin{InputList}
%
  \item[{\inp[nmlSEM::UseProfileForInitialization]{Use\_Profile\_For\_Initialization}}]
    This is a logical variable that, if set to \tru, will use the SEM
    profile to initialize the entire flow domain based on the $y$ coordinate
    value and the profile height variable.
    %
    This is very experimental and only really works for very simple geometries
    like a flat plate.
%
  \item[{\inp[nmlSEM::UseRandomSeedFromRestart]{Use\_Random\_Seed\_From\_Restart}}]
    As discussed in \sref{SEC:SEM:usage:restart}, when using an SEM boundary
    condition, the SEM data will also be written to the restart files including
    the seed used for the random number generator.\\[5pt]
    %
    This is a logical variable that, if \tru, will read the seed for
    the random number generator from the restart file and use it to reset the
    random number generator to exactly where it left off when the restart file
    was written, as if the simulation never stopped.\\[5pt]
    %
    The default value is \fls{} which will not use the restart
    information, and will instead use the current date and time to create a new
    \quoterm{random} seed for the random number generator.
    %
    This might have the tiniest effect on the \quoterm{randomness} of the
    strength and locations of the synthetic eddies in the SEM BC.
    %
    However, the total number of time steps needed for a simulation will almost
    certainly prevent this from having any kind of measurable effect on the
    simulation results.\\[5pt]
    %
    Very rarely (but it does happen), the SEM can randomly cause a simulation to
    die due to an NaN or other floating point exception.
    %
    Setting this to false is the only way to get past one of these occasions
    without modifying the solution through limiting \andor filtering.
%
%##%\FINALFIX{\newpage}%
  \item[{\inp[nmlSEM::UseSphericalEddyVolume]{Use\_Spherical\_Eddy\_Volume}}]
    This is a logical variable that determines if the volume of each synthetic
    eddy is treated as a box (\fls)[default] or as a sphere
    (\tru).
    %
    This input affects the number of eddies that will be used for
    an individual SEM boundary condition, since
    %
    \begin{equation*}
       [\textrm{\# of eddies}] =
       [\inptext{Box\_Num\_Eddies\_Multiplier}] \times
       \frac{[\textrm{Volume of SEM box}]}{[\textrm{Volume of synthetic eddy}]}
    \end{equation*}
%
  \item[{\inp[nmlSEM::IncludeThermoFluctuations]{Include\_Thermo\_Fluctuations}}]
    This is a logical variable that determines if fluctuations are added to the
    density and temperature profiles when computing the SEM boundary states.
    %
    This is has the value \fls{} by default meaning the density and
    temperature profiles are interpolated from the SEM profile with no
    fluctuations.
    %
    If this is set to \tru, fluctuations are added to the density and
    temperature profiles based on \linecite{Touber_2008}. Details on this are given
    in \sref{SEC:SEM:theory:ComputeBC}.
%
  \item[{\inp[nmlSEM::BoxNumEddiesMultiplier]{Box\_Num\_Eddies\_Multiplier}}]
    This is a floating point variable that is used as a multiplier when
    computing the total number of synthetic eddies
    used for the SEM boundary conditions, \thatis
    %
    \begin{equation*}
      [\textrm{\# of eddies}] = 
      [\inptext{Box\_Num\_Eddies\_Multiplier}] \times
      \frac{[\textrm{Volume of SEM box}]}{[\textrm{Volume of synthetic eddy}]}
    \end{equation*}
    %
    It should be noted that the number of eddies for an individual SEM BC are
    limited to be between a minimum of 5,000 and maximum of 500,000.
    %
    These limits prevent (1) too few eddies for being present to adequately
    model the flow, and (2) too many eddies that add little to the modeling
    but increase computational cost.
%
\end{InputList}

%###############################################################################

\subsubsubsection{Inputs Controlling Individual SEM Boundaries}[SEC:SEM:usage:semderivedtype]

\begin{nmldescitem}
&nml_bc
  ! Inputs for boundary group (1):
  BC_Input(1)%SEM%Profile_File        = ""
  BC_Input(1)%SEM%Is_Cylindrical      = F
  BC_Input(1)%SEM%Length_Scale        = 0.0  ! [m]
  BC_Input(1)%SEM%Convective_Velocity = 0.0  ! [m/s]
  BC_Input(1)%SEM%Profile_Height      = 0.0  ! [m]
  BC_Input(1)%SEM%Sub_In_BC_Grp       = 0
  BC_Input(1)%SEM%Sub_In_BC_Idx       = 0
/
\end{nmldescitem}

Inputs that are only applied to a specific SEM boundary are
specified in the \inp[nmlBC]{nml\_bc} input namelist block.
%
\begin{InputList}
%
  \item[{\BCinpBC[nmlBC::BCInput/SEM/ProfileFile]{SEM\%Profile\_File}}]
    This variable is a character string (max length of 150 characters) that
    specifies the path to the file containing the SEM profile used for the SEM
    boundary condition.
%
  \item[{\BCinpBC[nmlBC::BCInput/SEM/IsCylindrical]{SEM\%Is\_Cylindrical}}]
    This is a logical variable that is used to declare whether the SEM box is
    either rectangular or cylindrical. The default value of \fls{} indicates the
    SEM box is rectangular, whereas a value of \tru{} tells the code to use a
    cylindrical box for this SEM BC. See the last bullet in the list at the end
    of \sref{SEC:SEM:usage:profile:variables} for some additional technical
    details.
%
  \item[{\BCinpBC[nmlBC::BCInput/SEM/LengthScale]{SEM\%Length\_Scale}}]
    This is a floating point variable that gives the length scale (in meters) of
    the \quoterm{zone of influence} of each synthetic eddy. See
    \sref{SEC:SEM:usage:recommendations} for recommended values.
%
  \item[{\BCinpBC[nmlBC::BCInput/SEM/ConvectiveVelocity]%&
                 {SEM\%Convective\_Velocity}}]
    This is a floating point variable that gives the velocity (in m\slash{}s) at
    which the synthetic eddies are convected through the SEM box. See
    \sref{SEC:SEM:usage:recommendations} for recommended values.
%
\end{InputList}

The following three variables are also part of the \BCinptext{SEM}
structure, but they are only used internally within the code:
%
\begin{InputExample}
  \item[{\BCinpBC[nmlBC::BCInput/SEM/ProfileHeight]{SEM\%Profile\_Height}}]
  \item[{\BCinpBC[nmlBC::BCInput/SEM/SubInBCGrp]{SEM\%Sub\_In\_BC\_Grp}}]
  \item[{\BCinpBC[nmlBC::BCInput/SEM/SubInBCIdx]{SEM\%Sub\_In\_BC\_Idx}}]
\end{InputExample}
%
Defining any of these three variables in the input file will be ignored.
%
Their values will be determined after processing the SEM inputs, profile file,
and respective boundary faces.
%
This information will be written to the stdout display as part of the
startup process.
%
It can be used to verify that the profile data was
properly interpreted.

%###############################################################################

\subsubsubsection{Convective Velocity \slash{} Length Scale Recommendations}%
[SEC:SEM:usage:recommendations]
%
Previous testing seems to suggest a slower convective velocity should be used
when using a smaller length scale, and a higher convective velocity should be
used when using a larger length scale.
%
Let \uSEM{} be the value of
\BCinpBC[nmlBC::BCInput/SEM/ConvectiveVelocity]{SEM\%Convective\_Velocity}.
and let \lSEM{} be the value of
\BCinpBC[nmlBC::BCInput/SEM/LengthScale]{SEM\%Length\_Scale}.
%
%Additionally, let \uBC{} be the axial velocity at the top of the SEM profile
%(\thatis the edge velocity of the boundary layer in the SEM profile), and let
%\lBC{} be boundary layer height in the SEM profile.
%
Additionally, let \uBC{} be the axial velocity at the edge of the
boundary layer in the SEM profile file, and let
\lBC{} be boundary layer height in the SEM profile.
%
The following combinations all produce decent results in previous
tests of a channel flow:
%
\begin{itemize}[left=2.0em,align=left,style=nextline,noitemsep]
  \item[$\uSEM = \uBC/4\;$; $\quad\lSEM = \lBC/5$]
  \item[$\uSEM = \uBC/2\;$; $\quad\lSEM = \lBC/5$]
  \item[$\uSEM = \uBC/2\;$; $\quad\lSEM = \lBC/4$]
  \item[$\uSEM = \uBC/2\;$; $\quad\lSEM = \lBC/3.75$]
  \item[$\uSEM = \uBC/1\;$; $\quad\lSEM = \lBC/3.75$]
\end{itemize}
%
and are shown in \fref{FIG:SEM_Recommendations}.
%
From these results, recommended values for these two inputs are:
\\[8pt]
%\begin{tabular}{ @{\hspace{\descLabelLength}} rcccl @{} } 
\begin{tabular}{ @{\hspace{2.0em}} rcccl @{} } 
  $\uBC/4$ & $\!\!\!\le$ & $\!\!\!\uSEM$ & $\!\!\!\le$ & $\!\!\!\uBC$ \\
  $\lBC/5$ & $\!\!\!\le$ & $\!\!\!\lSEM$ & $\!\!\!\le$ & $\!\!\!\lBC/3.5$
\end{tabular}

\begin{figure}[h]
  \begin{centering}
    \insertimage[width=4.0in]{sem/sem_usem_lsem}
    \caption{Recommended SEM length scale and convective velocity.}
    \label{FIG:SEM_Recommendations}
  \end{centering}
\end{figure}

%###############################################################################

\subsubsubsection{Example SEM Boundary Condition}[SEC:SEM:usage:example]
 
\begin{nmldescitem}
&nml_bc
  ! Inputs for boundary group (1):
  BC_Input(1)%Set_BC_Value            = 5         ! 5 = SEM Inflow BC
  BC_Input(1)%Pressure                = 100000.0  ! [Pa]
  BC_Input(1)%Temperature             = 300.0     ! [K]
  BC_Input(1)%Mach                    = 0.2
  BC_Input(1)%SEM%Profile_File        = "[Path to profile file]"
  BC_Input(1)%SEM%Is_Cylindrical      = F
  BC_Input(1)%SEM%Length_Scale        = 0.0001    ! [m]
  BC_Input(1)%SEM%Convective_Velocity = 32.728    ! [m/s]
/
\end{nmldescitem}

In this example, if temperature is not specified in the SEM profile file, the
temperature for the entire SEM profile will be set to 300.0 K.

%%%%%%%%%%%%%%%%%%%%%%%%%%%%%%%%%%%%%%%%%%%%%%%%%%%%%%%%%%%%%%%%%%%%%%%%%%%%%%%%

Additionally, all face points with a wall distance larger than the \hSEM{}
determined from the profile file will be split off to a new subsonic inflow
boundary condition as discussed in \sref{SEC:SEM:usage:subsonic}.
%
The flow conditions for this new BC will be taken from (in this example) the
\inptext{Pressure}, \inptext{Temperature}, and \inptext{Mach} components of
\BCinpBC[nmlBC::BCInput][1]{}.

The \inptext{Convective\_Velocity} is chosen to provide
$U_{SEM}/U_{BL}$=0.5.

%\FIXME{\\
%Is this a working example?\\
%Was missing:
%  \inptext{Is\_Cylindrical},
%  \inptext{Length\_Scale}, and
%  \inptext{Convective\_Velocity}.
%}


%###############################################################################

